% SPDX-License-Identifier: Apache-2.0
\section{Every Name a \code{let} Can Want}
\label{sec:x-names}

A computed \code{let} in a step is a wire of the netlist, named as the
\code{let} is written. The netlist keeps one namespace for the unit's
fields, its ports and its wires, and Rust's rules do not reach into
it: a \code{let} may take a port's name, another \code{let}'s, a
field's, or a word VHDL or Verilog reserves, and every one of those
would have the netlist declare a name twice or use a word it cannot.

The lowering owns that namespace, so it decides. A wire takes the
\code{let}'s name where it is free, and \code{\_w}, \code{\_w2},
\code{\_w3} where it is not, which is one rule for all four cases. The
ports, the wires before it and the reserved words are known where the
attribute reads the step, so those are settled there; a field is not,
since \code{\#[lower]} reads the \code{impl} and never the struct, so
that part of the choice is a constant the compiler evaluates. A
constant is evaluated for every type the unit is lowered at, which is
why a generic unit needs no instantiation for the answer to be right.

None of it is silent. The netlist carries a line per renamed
\code{let} above the module, in the target's own comment syntax, so a
reader who meets a wire called \code{count\_w} is told which
\code{let} it came from:

\begin{lstlisting}[style=ascii]
// `let count` is the wire count_w.
// `let sum` is the wire sum_w.
// `let step` is the wire step_w.
// `let step` is the wire step_w2.
// `let buffer` is the wire buffer_w.
\end{lstlisting}

A \code{let} that only reads a register, a port or a number is an
alias rather than a wire, so it takes no name at all and may be called
anything. The unit below has one of those too, \code{same}, which
reads the register the field \code{count} holds.

One consequence is Rust's and not the netlist's: a \code{let} that
takes a port's name shadows the port in Rust as well, so the port is
driven above the \code{let} and read through the wire below it, as
\code{sum} is here.

\begin{figure}[htbp]
\centering
\input{names_timing.tex}
\caption{The counter: \code{count} and \code{last} are the fields,
\code{sum} the port the \code{let} shadows.}
\label{fig:x-names}
\end{figure}

\lstinputlisting[caption={\code{ex\_names.rs}. Run with \code{bazel run //lib/examples:ex\_names}.},label={lst:x-names}]{lib/examples/ex_names.rs}

\lstinputlisting[style=ascii,caption={What \code{ex\_names} printed when the build ran it: the counter, then the Verilog with the renamed wires and the lines that say which \code{let} each came from.},label={lst:x-names-out}]{out_names.txt}
